% Fibonacci-Pell nearest gaps: an all-exponent classification and
% quadratic-unit orbit rigidity
%
% Authors:
%   Dr. Denys Dutykh (Mathematics Department, Khalifa University of Science
%   and Technology, Abu Dhabi, UAE)
%   Prof. Laurent Vuillon (Univ. Savoie Mont Blanc, CNRS, LAMA, Chambery,
%   France)
%
\section{Why the Fibonacci and Pell arms}
\label{sec:two-arms}

A reader who has followed the proof may reasonably ask why this particular
pair of sequences carries the whole argument. The Markoff tree has
infinitely many branches; Fibonacci and Pell are two of them. This section
answers the question. The pairing is forced, not chosen: the mechanism used
throughout this paper exists on exactly these two branches, and the constant
$\gamma$ of \cref{eq:window-binet-phase} is the exchange rate between them.

Nothing in this section is used in the proofs of
\cref{thm:classification-intro,thm:orbit-intro,thm:local-clocks}. It is
included because it explains their shape, and because it delimits the
generalisations that can and cannot be attempted.

We are careful throughout to separate what is proved from what is only
verified on a finite family. Statements of the first kind are labelled as
propositions and theorems; statements of the second kind are labelled as
computational observations and are never used as hypotheses.

\subsection{Background from the Christoffel--Markoff correspondence}

We use the classical dictionary between Christoffel words and Markoff
numbers in the form given by Reutenauer \cite{Reutenauer2026}, whose origin is
Frobenius's observation of 1913 \cite{Frobenius1913} and Cohn's matrix
formulation. The one normalisation we need beyond the correspondence itself is
$m=\tfrac13\operatorname{tr}M_w=(M_w)_{12}$. In the transposed convention this
is \cite[Lemme~3.2]{Reutenauer2006}, which states that
$3N_{21}=\operatorname{tr}N$ for the corresponding conjugate; the supplement
also verifies the identity directly, in the convention used here, on every
Christoffel word of depth at most eleven. Let

\begin{equation}
 \mathsf A=\begin{pmatrix}2&1\\1&1\end{pmatrix},
 \qquad
 \mathsf B=\begin{pmatrix}5&2\\2&1\end{pmatrix},
 \label{eq:cohn-generators}
\end{equation}

and for a word $w$ over $\{a,b\}$ write $M_w$ for the corresponding product,
$a\mapsto\mathsf A$ and $b\mapsto\mathsf B$. If $w$ is a Christoffel word,
then $M_w\in\mathrm{SL}_2(\Z)$ has positive entries and

\begin{equation}
 M_w=\begin{pmatrix}A&m\\C&D\end{pmatrix},
 \qquad
 A+D=3m,
 \qquad
 1\leqslant D\leqslant m,
 \label{eq:cohn-normalisation}
\end{equation}

where $m$ is the Markoff number attached to $w$. Two consequences are used
below. First, $AD-mC=1$ and $A\equiv-D\pmod m$ give $D^2\equiv-1\pmod m$, so
the residue

\begin{equation}
 u=A-2m=m-D
 \label{eq:characteristic-root}
\end{equation}

lies in $[0,m)$ and satisfies $u^2+1\equiv0\pmod m$; it is the
characteristic root of the Markoff number. Second, the Markoff word attached
to $m$, that is the Christoffel word of slope $u/(m-u)$, has length exactly
$m$, with

\begin{equation}
 m_b=u\ \text{letters }b,
 \qquad
 m_a=m-u\ \text{letters }a.
 \label{eq:letter-counts}
\end{equation}

We write $\varphi=(1+\sqrt5)/2$ throughout, as elsewhere in this paper.

\subsection{The writing and the letter counts determine each other}

The seed data of this paper is a primitive two-square writing $m=p^2+q^2$,
whereas the word side supplies the letter counts \cref{eq:letter-counts}. The
passage between the two is the Gauss correspondence between primitive writings
and square roots of $-1$, and it takes the concrete form of a congruence with
an explicit quotient. We record it because it turns the first row of
\cref{tab:dictionary} into an identity rather than an analogy, and because the
orientation ambiguity it carries is exactly the ambiguity
$u\leftrightarrow m-u$ of the characteristic root.

\begin{proposition}[Gauss congruence and its orientation]
\label{prop:gauss-congruence}
Let $p$ and $q$ be coprime positive integers, put $m=p^2+q^2$, and let
$u\in[0,m)$ be the residue determined by $uq\equiv p\pmod m$. Then

\begin{enumerate}[label=\textup{(\roman*)}]
 \item $u^2+1\equiv0\pmod m$;
 \item $q-p+u(p+q)\equiv0\pmod m$, and the quotient

 \begin{equation}
  k=\frac{u(p+q)+q-p}{m}
  \label{eq:gauss-quotient}
 \end{equation}

 is the unique integer of the open interval $(0,p+q)$ satisfying
 $kq\equiv1\pmod{p+q}$;
 \item exchanging $p$ and $q$ replaces $u$ by $m-u$ and $k$ by $p+q-k$; the
 companion congruence is $p-q+(m-u)(p+q)\equiv0\pmod m$, and its quotient
 $p+q-k$ is the inverse of $p$ modulo $p+q$.
\end{enumerate}
\end{proposition}

\begin{proof}
\emph{The residue is well defined.} A common divisor of $q$ and $m=p^2+q^2$
divides $p^2$, hence divides $\gcd(p,q)^2=1$. So $q$ is invertible modulo $m$,
and $u=pq^{-1}\bmod m$ exists and is unique in $[0,m)$.

\emph{Part (i).} Reducing $m=p^2+q^2$ modulo $m$ gives $p^2\equiv-q^2$.
Multiplying by the square of $q^{-1}$,

\begin{equation}
 u^2=p^2\bigl(q^{-1}\bigr)^2\equiv-q^2\bigl(q^{-1}\bigr)^2=-1\pmod m.
\end{equation}

\emph{Part (ii), the congruence.} Expanding and using $p^2\equiv-q^2$ once,

\begin{equation}
 u(p+q)\equiv pq^{-1}(p+q)=p+p^2q^{-1}\equiv p-q^2q^{-1}=p-q\pmod m,
 \label{eq:gauss-congruence-core}
\end{equation}

which is the assertion.

\emph{Part (ii), the quotient lies in $(0,p+q)$.} First $u\neq0$: otherwise
$p\equiv0\pmod m$, which is impossible because $1\leqslant p<m$. Hence
$1\leqslant u\leqslant m-1$, and the numerator of \cref{eq:gauss-quotient} is
bounded on both sides by

\begin{equation}
 2q=(p+q)+q-p\leqslant u(p+q)+q-p\leqslant(m-1)(p+q)+q-p=m(p+q)-2p.
\end{equation}

The left end is positive and the right end is smaller than $m(p+q)$, so
dividing by $m$ gives $0<k<p+q$; and $k$ is an integer because the numerator is
divisible by $m$ by the previous step.

\emph{Part (ii), the quotient is an inverse.} By the definition of $u$ there is
an integer $a$ with $uq=p+am$. Multiplying $km=u(p+q)+q-p$ by $q$ and
substituting,

\begin{equation}
 kmq=uq(p+q)+q^2-pq=(p+am)(p+q)+q^2-pq=p^2+q^2+am(p+q),
\end{equation}

the term $pq$ cancelling. The right-hand side is $m+am(p+q)$, so dividing
through by $m$,

\begin{equation}
 kq=1+a(p+q).
 \label{eq:gauss-quotient-inverse}
\end{equation}

Thus $kq\equiv1\pmod{p+q}$. Since $\gcd(q,p+q)=\gcd(q,p)=1$, that inverse is
unique modulo $p+q$, and by the previous step $k$ is its representative in
$(0,p+q)$.

\emph{Part (iii).} Exchanging $p$ and $q$ replaces $u$ by the residue $u'$
determined by $u'p\equiv q$, that is $u'=qp^{-1}\bmod m$. Then $uu'\equiv1$,
and multiplying $u^2\equiv-1$ by $u'$ gives $u\equiv u^2u'\equiv-u'$; since
both lie in $[0,m)$ and $u\neq0$, this means $u'=m-u$. Parts (i) and (ii)
applied to the pair $(q,p)$ give the companion congruence together with the
statement that its quotient $k'=\bigl(u'(p+q)+p-q\bigr)/m$ is the inverse of
$p$ modulo $p+q$. Finally, adding the two numerators and using $u+u'=m$,

\begin{equation}
 (k+k')m=(u+u')(p+q)=m(p+q),
\end{equation}

so $k'=p+q-k$.
\end{proof}

\begin{proposition}[The writing attached to a root of $-1$]
\label{prop:writing-from-root}
Let $m\geqslant5$ and let $u$ be an integer with $u^2+1\equiv0\pmod m$. Then
there is exactly one pair $(p,q)$ of positive integers with $m=p^2+q^2$ and
$p\equiv uq\pmod m$, and that pair is coprime.
\end{proposition}

\begin{proof}
\emph{A small solution of the congruence.} Put $n=\lfloor\sqrt m\rfloor$ and
consider the $(n+1)^2$ pairs $(s,t)$ of integers with $0\leqslant s,t\leqslant n$.
Since $(n+1)^2>m$, two distinct such pairs $(s_1,t_1)$ and $(s_2,t_2)$ satisfy
$s_1-ut_1\equiv s_2-ut_2\pmod m$. Put $x=s_1-s_2$ and $y=t_1-t_2$, so that

\begin{equation}
 x\equiv uy\pmod m,
 \qquad
 \abs x\leqslant n,
 \qquad
 \abs y\leqslant n,
 \label{eq:thue-pair}
\end{equation}

with $(x,y)\neq(0,0)$. Neither coordinate vanishes. If $y=0$ then $m\mid x$
while $0<\abs x\leqslant n<m$; and if $x=0$ then $m\mid uy$, hence $m\mid y$
because $\gcd(u,m)=1$, while $0<\abs y\leqslant n<m$.

\emph{The pair represents $m$.} Squaring \cref{eq:thue-pair} and using
$u^2\equiv-1$,

\begin{equation}
 x^2+y^2\equiv u^2y^2+y^2=(u^2+1)y^2\equiv0\pmod m,
\end{equation}

while $0<x^2+y^2\leqslant2n^2\leqslant2m$. So $x^2+y^2$ is either $m$ or $2m$.
Suppose it were $2m$. Then $\abs x=\abs y=n$ and $m=n^2$, so $x=\delta n$ and
$y=\delta'n$ for two signs $\delta,\delta'\in\{1,-1\}$, and
$x-uy=n(\delta-u\delta')$. Since $m=n^2$ divides $x-uy$, we get
$n\mid\delta-u\delta'$, hence $u\equiv\delta\delta'\pmod n$ and therefore
$u^2\equiv1\pmod n$. But $n\mid m$ and $m\mid u^2+1$ give $u^2\equiv-1\pmod n$,
so $n\mid2$ and $m=n^2\leqslant4$, contradicting $m\geqslant5$. Hence
$x^2+y^2=m$.

\emph{The pair is coprime.} Let $g=\gcd(x,y)$ and write $x=gx_1$, $y=gy_1$.
Then $g^2\mid x^2+y^2=m$; put $m_1=m/g^2=x_1^2+y_1^2$. The divisibility
$m\mid x-uy$ reads $g^2m_1\mid g(x_1-uy_1)$, that is

\begin{equation}
 x_1\equiv uy_1\pmod{gm_1}.
\end{equation}

Squaring it and using $u^2\equiv-1$ modulo $g^2m_1$, hence a fortiori modulo
$gm_1$,

\begin{equation}
 m_1=x_1^2+y_1^2\equiv(u^2+1)y_1^2\equiv0\pmod{gm_1}.
\end{equation}

A positive integer divisible by $gm_1$ and equal to $m_1$ forces $g=1$.

\emph{Normalisation.} Changing $(x,y)$ into $(-x,-y)$ leaves
\cref{eq:thue-pair} unchanged, so we may assume $y>0$. If $x>0$, take
$(p,q)=(x,y)$. If $x<0$, put $p=y$ and $q=-x$, both positive; multiplying
$x\equiv uy$ by $u$ and using $u^2\equiv-1$ gives $ux\equiv-y$, that is
$y\equiv u(-x)$, which is $p\equiv uq$. In both cases $m=p^2+q^2$ with
$\gcd(p,q)=1$ and $p\equiv uq\pmod m$.

\emph{Uniqueness.} Let $(p,q)$ and $(P,Q)$ be pairs of positive integers with
$m=p^2+q^2=P^2+Q^2$, $p\equiv uq$ and $P\equiv uQ$ modulo $m$. Then

\begin{equation}
 pQ-qP\equiv uqQ-quQ=0\pmod m.
\end{equation}

On the other hand $pQ-qP$ is the inner product of $(p,q)$ with $(Q,-P)$, so by
the Cauchy--Schwarz inequality

\begin{equation}
 \abs{pQ-qP}\leqslant\sqrt{p^2+q^2}\,\sqrt{P^2+Q^2}=m,
\end{equation}

with equality only if $(p,q)$ and $(Q,-P)$ are proportional. That is
impossible: $p=\mu Q$ would need $\mu>0$ and $q=-\mu P$ would need $\mu<0$.
Hence $\abs{pQ-qP}<m$, and the congruence forces $pQ=qP$. Since
$\gcd(p,q)=1$, this gives $P=\kappa p$ and $Q=\kappa q$ for a positive integer
$\kappa$, and then $m=P^2+Q^2=\kappa^2m$ gives $\kappa=1$.
\end{proof}

\begin{corollary}[The writing and the letter counts determine each other]
\label{cor:writing-letter-counts}
Let $m\geqslant5$ be a Markoff number and $u$ its characteristic root
\cref{eq:characteristic-root}. Then $m$ has exactly one writing $m=p^2+q^2$ in
positive integers with $p\equiv uq\pmod m$, that writing is primitive, and the
letter counts \cref{eq:letter-counts} of the Markoff word of $m$ satisfy

\begin{equation}
 m_b(p+q)\equiv p-q,
 \qquad
 m_a(p+q)\equiv q-p
 \pmod m.
 \label{eq:letter-counts-congruence}
\end{equation}

Exchanging $p$ with $q$ replaces $u$ by $m-u$ and exchanges $m_a$ with $m_b$.
\end{corollary}

\begin{proof}
The characteristic root satisfies $u^2+1\equiv0\pmod m$ by
\cref{eq:characteristic-root}, so \cref{prop:writing-from-root} supplies the
writing, its primitivity and its uniqueness. Its defining congruence
$p\equiv uq$ says that the residue attached to $(p,q)$ in
\cref{prop:gauss-congruence} is $u$ itself, so
\cref{eq:gauss-congruence-core} gives $u(p+q)\equiv p-q$; by
\cref{eq:letter-counts} that residue is the letter count $m_b$, which is the
first congruence. The second follows from $m_a=m-m_b$, since then
$m_a(p+q)\equiv-u(p+q)\equiv q-p$. The last assertion is part (iii) of
\cref{prop:gauss-congruence}.
\end{proof}

For a Markoff number the two-square writing and the Markoff word are therefore
two readings of one object, and each is computed from the other: the writing
from the characteristic root by \cref{prop:writing-from-root}, the letter
counts from the writing by \cref{eq:letter-counts-congruence}. The only
ambiguity in the passage is the exchange $u\leftrightarrow m-u$, which is the
choice of an ordering of the pair $(p,q)$. The next subsection computes both
sides explicitly on the two extreme arms of the tree, where the letter counts
turn out to be the coordinates of a power of the fundamental unit.

\subsection{Letter counts on the two arms are unit coordinates}

The leftmost branch of the Christoffel tree consists of the words $a^kb$
and produces the Fibonacci Markoff numbers; the rightmost consists of the
words $ab^k$ and produces the Pell Markoff numbers. On both, the letter
counts are not merely reminiscent of the fundamental unit: they are its
coordinates.

\begin{proposition}[Letter counts on the two arms]
\label{prop:arm-letter-counts}
For $k\geqslant1$,

\begin{equation}
 (m_a,m_b)=(F_{2k},F_{2k-1})
 \quad\text{when}\quad
 m=F_{2k+1},
 \label{eq:fib-arm-counts}
\end{equation}

and

\begin{equation}
 (m_a,m_b)=(C_{2k},P_{2k})
 \quad\text{when}\quad
 m=P_{2k+1}.
 \label{eq:pell-arm-counts}
\end{equation}

Equivalently, $(m_a,m_b)$ is the coordinate vector of $\varphi^{2k}$ in the
basis $(\varphi,1)$ of $\Z[\varphi]$ in the first case, and of $\lambda^{2k}$
in the basis $(1,\sqrt2)$ of $\Z[\sqrt2]$ in the second.
\end{proposition}

\begin{proof}
Both cases are a direct computation with the Cohn generators
\cref{eq:cohn-generators}, using \cref{eq:characteristic-root} in the form
$m=(M_w)_{12}$ and $u=(M_w)_{11}-2m$. Two elementary matrix identities drive
them. Writing

\begin{equation}
 \mathsf F=\begin{pmatrix}1&1\\1&0\end{pmatrix},
 \qquad
 \mathsf P=\begin{pmatrix}2&1\\1&0\end{pmatrix},
 \label{eq:fib-pell-matrices}
\end{equation}

one checks at once that $\mathsf F^2=\mathsf A$ and
$\mathsf P^2=\mathsf B$, and the standard induction gives

\begin{equation}
 \mathsf F^{\,r}=\begin{pmatrix}F_{r+1}&F_r\\F_r&F_{r-1}\end{pmatrix},
 \qquad
 \mathsf P^{\,r}=\begin{pmatrix}P_{r+1}&P_r\\P_r&P_{r-1}\end{pmatrix}.
 \label{eq:power-matrices}
\end{equation}

\emph{Fibonacci arm.} The word is $w=a^{k-1}b$, so
$M_w=\mathsf A^{k-1}\mathsf B=\mathsf F^{2k-2}\mathsf B$. Multiplying
\cref{eq:power-matrices} by $\mathsf B$ gives

\begin{equation}
 M_w=
 \begin{pmatrix}
 5F_{2k-1}+2F_{2k-2} & 2F_{2k-1}+F_{2k-2}\\
 5F_{2k-2}+2F_{2k-3} & 2F_{2k-2}+F_{2k-3}
 \end{pmatrix}.
\end{equation}

Its $(1,2)$ entry is
$2F_{2k-1}+F_{2k-2}=F_{2k-1}+(F_{2k-1}+F_{2k-2})=F_{2k-1}+F_{2k}=F_{2k+1}$,
so $m=F_{2k+1}$. Subtracting twice this from the $(1,1)$ entry,

\begin{equation}
 u=(5F_{2k-1}+2F_{2k-2})-2(2F_{2k-1}+F_{2k-2})=F_{2k-1}.
\end{equation}

Hence $m_b=F_{2k-1}$ and $m_a=F_{2k+1}-F_{2k-1}=F_{2k}$, which is
\cref{eq:fib-arm-counts}. The coordinate statement is Binet's identity
$\varphi^{2k}=F_{2k}\varphi+F_{2k-1}$, immediate by induction from
$\varphi^2=\varphi+1$.

\emph{Pell arm.} The word is $w=ab^{k}$, so
$M_w=\mathsf A\mathsf B^{k}=\mathsf A\mathsf P^{2k}$, and
\cref{eq:power-matrices} gives

\begin{equation}
 M_w=
 \begin{pmatrix}
 2P_{2k+1}+P_{2k} & 2P_{2k}+P_{2k-1}\\
 P_{2k+1}+P_{2k} & P_{2k}+P_{2k-1}
 \end{pmatrix}.
\end{equation}

Its $(1,2)$ entry is $2P_{2k}+P_{2k-1}=P_{2k+1}$ by the Pell recursion, so
$m=P_{2k+1}$. Subtracting twice this from the $(1,1)$ entry,

\begin{equation}
 u=(2P_{2k+1}+P_{2k})-2P_{2k+1}=P_{2k}.
\end{equation}

Hence $m_b=P_{2k}$ and $m_a=P_{2k+1}-P_{2k}=C_{2k}$, the last equality being
the recursion $P_{r+1}=C_r+P_r$ recorded after
\cref{eq:pell-coordinates}. This is \cref{eq:pell-arm-counts}.

Both computed roots do satisfy $u^2+1\equiv0$ modulo $m$, as
\cref{eq:characteristic-root} requires, and in both cases this is a
determinant. Since $\det\mathsf F=\det\mathsf P=-1$, taking determinants in
\cref{eq:power-matrices} at the even exponent $r=2k$ gives

\begin{equation}
 F_{2k+1}F_{2k-1}-F_{2k}^2=1,
 \qquad
 P_{2k+1}P_{2k-1}-P_{2k}^2=1.
 \label{eq:cassini-even}
\end{equation}

The second congruence is immediate: $P_{2k}^2\equiv-1\pmod{P_{2k+1}}$. The
first gives $F_{2k}^2\equiv-1\pmod{F_{2k+1}}$, and since
$F_{2k-1}=F_{2k+1}-F_{2k}\equiv-F_{2k}$, also
$F_{2k-1}^2\equiv F_{2k}^2\equiv-1\pmod{F_{2k+1}}$.
\end{proof}

\begin{remark}[A numerical coincidence explained]
\label{rem:377-233}
The case $k=7$ of \cref{eq:fib-arm-counts} reads $m=F_{15}=610$ and
$(m_a,m_b)=(377,233)$, and the corresponding unit identity is
$\varphi^{14}=377\varphi+233$, which is \cref{eq:phi14}. The two occurrences
of the pair $(377,233)$ in this paper, one as a threshold in the
even-branch estimate and one as the letter counts of the Markoff word of
$610$, are therefore the same identity read twice. The first nontrivial Pell
case is $\lambda^4=17+12\sqrt2$: the Markoff word of $29$ has $17$ letters
$a$ and $12$ letters $b$.

The passage from the two-square writing to the letter counts is not an
analogy to be tested on examples: \cref{cor:writing-letter-counts} makes it an
identity, proved for every Markoff number. It can also be watched on the word
side, where it takes a purely combinatorial form. Starting from the two-square
pair one recovers the half directive word by iterated palindromic closure and
doubles it by reversal complement. At $m=29$ the pair $(2,5)$ returns the half
directive word $aab$, whose reversal complement $abb$ prefixes it into the
directive word $abbaab$, and the Markoff word
it builds has the $17$ letters $a$ and $12$ letters $b$ that
\cref{eq:letter-counts} predicts from the characteristic root $u=12$, a
quantity that mentions neither the pair nor any palindrome.
\end{remark}

\subsection{Every branch carries a quadratic Perron unit}

Fix a Christoffel word $u_0$ with Markoff number $m_0$, and follow the branch
of the tree that keeps $m_0$ fixed, so that the node words are
$u_0^{\,k}v_0$ for a fixed $v_0$. Write $m_k$ for the Markoff number of the
$k$-th node and $u_k$ for its characteristic root.

\begin{proposition}[Branch recurrence]
\label{prop:branch-recurrence}
The vector $(m_a,m_b)_k=(m_k-u_k,u_k)$ satisfies

\begin{equation}
 (m_a,m_b)_{k+1}
 =3m_0\,(m_a,m_b)_k-(m_a,m_b)_{k-1}.
 \label{eq:branch-recurrence}
\end{equation}

Its characteristic polynomial is $x^2-3m_0x+1$, whose dominant root

\begin{equation}
 \beta(m_0)=\frac{3m_0+\sqrt{9m_0^2-4}}2
 \label{eq:branch-perron}
\end{equation}

is a quadratic unit, hence a Perron number, of the order of discriminant
$9m_0^2-4$.
\end{proposition}

\begin{proof}
Write $P_k=M_{u_0^{\,k}v_0}$ for the Cohn matrix of the $k$-th node. Since
the Cohn correspondence is multiplicative on concatenation,
$P_k=M_{u_0}^{\,k}M_{v_0}$.

\emph{Step 1: the matrix recurrence.} The matrix $M_{u_0}$ lies in
$\mathrm{SL}_2(\Z)$ and has trace $3m_0$ by
\cref{eq:cohn-normalisation}. Its characteristic polynomial is therefore
$x^2-3m_0x+1$, and Cayley--Hamilton gives
$M_{u_0}^2=3m_0M_{u_0}-I$. Multiplying on the right by
$M_{u_0}^{\,k-1}M_{v_0}$ yields

\begin{equation}
 P_{k+1}=3m_0P_k-P_{k-1},
 \label{eq:matrix-recurrence}
\end{equation}

an identity between matrices, hence between each of their four entries
separately.

\emph{Step 2: the two entries we need.} By \cref{eq:cohn-normalisation}, the
Markoff number is the $(1,2)$ entry, $m_k=(P_k)_{12}$, and the $(1,1)$ entry
is $(P_k)_{11}$. Both obey the scalar recurrence
$x_{k+1}=3m_0x_k-x_{k-1}$, by Step 1.

\emph{Step 3: the characteristic root obeys it too.} By
\cref{eq:characteristic-root}, $u_k=(P_k)_{11}-2m_k$. A fixed integer linear
combination of two solutions of a linear recurrence is again a solution, so
$u_k$ satisfies $x_{k+1}=3m_0x_k-x_{k-1}$. The constancy of the coefficient
$2$ is exactly the normalisation $1\leqslant D\leqslant m$ recorded in
\cref{eq:cohn-normalisation}, which holds at every node.

\emph{Step 4: assembling the vector.} Both coordinates of
$(m_a,m_b)_k=(m_k-u_k,u_k)$ are integer linear combinations of $m_k$ and
$u_k$, so both satisfy the same recurrence, which is
\cref{eq:branch-recurrence}. Its characteristic polynomial is
$x^2-3m_0x+1$; since $m_0\geqslant1$ the discriminant $9m_0^2-4$ is positive
and not a square for $m_0\geqslant2$, and the dominant root is
\cref{eq:branch-perron}. It is a unit because the product of the two roots is
the constant term $1$.
\end{proof}

The two arms are the cases $m_0=1$ and $m_0=2$, where
$\beta(1)=(3+\sqrt5)/2=\varphi^2$ and
$\beta(2)=3+2\sqrt2=\lambda^2$.

\begin{remark}[The order in which $\beta(m_0)$ lives is classical]
\label{rem:frobenius-conic}
The discriminant $9m_0^2-4$ of \cref{eq:branch-perron} is not new to Markoff
theory. At a fixed largest entry $c$, the Markoff equation can be written as
the pair of identities

\begin{equation}
 (a-b)^2+c^2=ab(3c-2),
 \qquad
 (a+b)^2+c^2=ab(3c+2),
 \label{eq:frobenius-identities}
\end{equation}

which are printed in Frobenius \cite[p.~601 of the reprint]{Frobenius1913}
and are quoted as equations~(3) and~(4) by Zhang \cite{Zhang2007}, and whose
cross product is a conic of discriminant $(3c-2)(3c+2)=9c^2-4$. That conic is
itself classical: in the halved coordinates $A=(a+b)/2$ and $B=(a-b)/2$ it is
printed as

\begin{equation}
 (2-3c)A^2+(2+3c)B^2=-c^2
 \label{eq:srinivasan-conic}
\end{equation}

in \cite[eq.~(2.3)]{Srinivasan2009}. So the
order that carries the branch unit $\beta(m_0)$ is exactly the order attached
to $m_0$ by the classical theory, and the two extreme arms of
\cref{thm:half-step-criterion} are the two branches whose unit happens to be
a square in it. They are also the source of the sharpest known unconditional
uniqueness criteria, which run modulo $3c\pm2$ rather than modulo $c$:
\cite[Theorem~2]{Zhang2007} settles every Markoff number for which one of
$3c-2$, $3c+2$ is a prime power, four times a prime power, or eight times a
prime power, and \cite[Theorem~2.3]{Srinivasan2009} settles every Markoff
number whose $3c-2$ or $3c+2$ has a prime power for greatest odd divisor. The
two hypotheses coincide on Markoff numbers, because $3c\pm2$ is odd for odd
$c$, while for even $c$ the quotients $(3c-2)/4$ and $(3c+2)/8$ are odd
integers \cite[Lemma~4.2]{Srinivasan2009}.

Srinivasan works in exactly the order met here. Her uniqueness criterion
\cite[Theorem~5.4]{Srinivasan2009} asks that every ambiguous form lying in the
principal genus of the class group of discriminant $9c^2-4$ correspond to a
divisor of $3c-2$, and for a $c$ that is not a prime power she proves that
Markoff uniqueness at $c$ is \emph{equivalent} to a statement about
$\Q(\sqrt{9c^2-4})$ \cite[Theorem~5.3]{Srinivasan2009}. The order attached to
a branch is therefore not an artefact of our construction; it is where the
difficulty of the conjecture is already known to sit.

One consequence is worth recording, because it explains why
\cref{thm:half-step-criterion} is stated for a unit of an arbitrary real
quadratic order rather than of this one. The branch unit $\beta(m_0)$ has norm
$1$ in the order of discriminant $9m_0^2-4$, and a square root of it of norm
$-1$ need not lie in that order. At $m_0=1$ it does: the discriminant is $5$,
the order is the maximal order of $\Q(\sqrt5)$, and $\varphi$ lies in it. At
$m_0=2$ it does not: the discriminant is $32$, the order is $\Z[2\sqrt2]$, and
$\lambda=1+\sqrt2$ lies instead in the maximal order $\Z[\sqrt2]$, of
discriminant $8$. Restricting $\varepsilon$ to the order of $\beta(m_0)$ would
therefore have discarded the Pell arm. We use none of this, and record it only
because it locates \cref{prop:branch-recurrence} inside the classical
picture.
\end{remark}

\subsection{The mechanism exists on exactly two branches}

\Cref{prop:branch-recurrence} says that a quadratic unit governs every
branch. That is not what the proofs in this paper consume. What they consume
is the semiconjugacy of \cref{prop:seed-semiconjugacy}: the map
$\mathsf S$ advances the two-square data by one step, and $\mathsf S^2$
realises $\beta$. In other words the argument needs a \emph{square root} of
$\beta(m_0)$, acting integrally. The following theorem shows that this is a
severe restriction.

\begin{theorem}[Trace criterion for a half-step]
\label{thm:half-step-criterion}
Let $m_0$ be a Markoff number and suppose $\beta(m_0)=\varepsilon^2$ for a
unit $\varepsilon$ of a real quadratic order. Then

\begin{equation}
 3m_0-2\ \text{is a perfect square},
 \label{eq:half-step-criterion}
\end{equation}

and $\varepsilon$ has norm $-1$ and trace $\sqrt{3m_0-2}$.
\end{theorem}

\begin{proof}
Let $t=\varepsilon+\overline\varepsilon$ be the trace of $\varepsilon$ and
$N(\varepsilon)=\varepsilon\overline\varepsilon=\pm1$ its norm; both are
rational integers. Squaring,

\begin{equation}
 \beta(m_0)+\overline{\beta(m_0)}
 =\varepsilon^2+\overline\varepsilon^{\,2}
 =t^2-2N(\varepsilon).
\end{equation}

The left-hand side is the trace of $\beta(m_0)$, which is $3m_0$ by
\cref{eq:branch-perron}. Hence

\begin{equation}
 t^2=3m_0+2N(\varepsilon).
\end{equation}

Two cases arise. If $N(\varepsilon)=1$ then $t^2=3m_0+2$. But every square is
congruent to $0$ or $1$ modulo $3$, while $3m_0+2\equiv2\pmod3$, so this case
is impossible. Therefore $N(\varepsilon)=-1$ and $t^2=3m_0-2$, which is
\cref{eq:half-step-criterion}.
\end{proof}

For the two arms the criterion is satisfied with room to spare:
$3\cdot1-2=1$ gives $t=1$ and $\varepsilon=\varphi$, while
$3\cdot2-2=4$ gives $t=2$ and $\varepsilon=\lambda$. Both have norm $-1$, as
the theorem requires, and both are the units of this paper.

\begin{remark}[Computational scope]
\label{rem:half-step-scan}
Among all $3502$ Markoff numbers below $10^{60}$, and again among the
$4097$ produced by the Christoffel tree to depth $11$, the criterion
\cref{eq:half-step-criterion} is met by exactly three: $m_0=1$, $m_0=2$ and
$m_0=34$, with $3\cdot34-2=100$ and $\varepsilon=5+\sqrt{26}$. This is a
verification on a finite family and not a theorem: we do not know whether
$3m-2$ is a perfect square for infinitely many Markoff numbers $m$, and
nothing below depends on the answer. The certificate is
\texttt{check\_two\_arms.py} in the supplement.
\end{remark}

The third candidate is excluded, and here the finiteness is genuine, because
the statement concerns one explicit linear system.

\begin{proposition}[No half-step at $m_0=34$]
\label{prop:no-half-step-34}
The branch fixing $m_0=34$ carries no integral half-step on the two-square
data.
\end{proposition}

\begin{proof}
The successive two-square writings on that branch are
$(5,8)$, $(34,89)$, $(513,814)$ and $(3468,9077)$, attached to the Markoff
numbers $89$, $9077$, $925765$ and $94418953$. A half-step would be a matrix
$\mathsf T\in\mathrm{GL}_2(\Z)$ carrying each writing to the next. The first
two conditions already determine $\mathsf T$ over $\Q$, because $(5,8)$ and
$(34,89)$ are linearly independent, and solving the resulting linear system
gives

\begin{equation}
 \mathsf T=\frac1{173}
 \begin{pmatrix}-1078&1409\\1409&1044\end{pmatrix},
\end{equation}

whose entries are not integers. Moreover this $\mathsf T$ sends
$(513,814)$ to

\begin{equation}
 \left(\frac{593912}{173},\ \frac{1572633}{173}\right)
 \ne(3468,9077),
\end{equation}

so no matrix at all, integral or not, satisfies the first three conditions
simultaneously.
\end{proof}

Thus the criterion admits three branches, and only two of them survive. The
pairing of $\varphi$ against $\lambda$ in this paper is therefore not a
choice among infinitely many; it is the complete list of branches on which
the semiconjugacy of \cref{prop:seed-semiconjugacy} can exist at all.

\subsection{The two arms are the two constant directive words}

There is a second, purely combinatorial way to see the same pair, and it
explains where $\gamma$ comes from.

A Markoff directive word has all its partial quotients in $\{1,2\}$. There
are exactly two constant such sequences, and their values are

\begin{equation}
 [1;1,1,1,\ldots]=\varphi,
 \qquad
 [2;2,2,2,\ldots]=1+\sqrt2=\lambda.
 \label{eq:constant-directives}
\end{equation}

These are the Fibonacci and Pell arms. The next proposition converts this
into a statement about the cost of a directive word, that is, the number of
steps of the subtractive Euclidean algorithm on the two-square pair, which
equals the sum of the partial quotients minus one.

\begin{proposition}[Arm costs and the meaning of $\gamma$]
\label{prop:arm-costs}
Let $\ell(p,q)$ denote the cost of the pair $(p,q)$. On the Fibonacci arm,
where $m=F_{2k+1}$ and $(p,q)=(F_k,F_{k+1})$,

\begin{equation}
 \ell(F_k,F_{k+1})=k-1;
\end{equation}

on the Pell arm, where $m=P_{2k+1}$ and $(p,q)=(P_k,P_{k+1})$,

\begin{equation}
 \ell(P_k,P_{k+1})=2k-1.
\end{equation}

Consequently, comparing the two arms at equal magnitude,

\begin{equation}
 \frac{\ell\ \text{on the Pell arm}}{\ell\ \text{on the Fibonacci arm}}
 \longrightarrow
 \frac{2\log\varphi}{\log\lambda}
 =2\gamma
 =1.09195\ldots
 \label{eq:cost-ratio}
\end{equation}
\end{proposition}

\begin{proof}
The continued fraction of $F_{k+1}/F_k$ is $[1,1,\ldots,1,2]$ with $k-1$
partial quotients, by induction from $F_{k+1}=F_k+F_{k-1}$; its quotient sum
is $(k-2)+2=k$, so the cost is $k-1$. The continued fraction of
$P_{k+1}/P_k$ is $[2,2,\ldots,2]$ with $k$ partial quotients, by induction
from $P_{k+1}=2P_k+P_{k-1}$; its quotient sum is $2k$, so the cost is
$2k-1$.

For the ratio, Binet's formulas give

\begin{equation}
 F_{2k+1}=\frac{\varphi^{2k+1}}{\sqrt5}+O(1),
 \qquad
 P_{2k+1}=\frac{\lambda^{2k+1}}{2\sqrt2}+O(1).
\end{equation}

Solving each for its index at a
common magnitude $m$ gives $k=\log m/(2\log\varphi)+O(1)$ on the Fibonacci
arm and $k=\log m/(2\log\lambda)+O(1)$ on the Pell arm. Substituting into
the two cost formulas gives $\log m/(2\log\varphi)+O(1)$ and
$\log m/\log\lambda+O(1)$ respectively, whose ratio tends to
$2\log\varphi/\log\lambda$.
\end{proof}

So $\gamma$ is not a computational artefact of the reduction. It is the
asymptotic exchange rate between the directive-word costs of the two arms,
and the rational $p/Q$ of \cref{eq:common-convergent} is a certified
rational approximation to that exchange rate.

\begin{remark}[What is not claimed]
\label{rem:not-extremal}
It is tempting to strengthen \cref{eq:constant-directives} into a statement
that the two arms are the extremes of the growth rate of admissible directive
words. Only one half of that is true. Among all compositions of a given
quotient sum $N$, the continuant is maximal at $(1,1,\ldots,1)$ and equals
$F_{N+1}$, so the Fibonacci arm does grow fastest and
$\ell(p,q)\geqslant\log(p+q)/\log\varphi+O(1)$ for every writing. The Pell
arm, however, is not the slowest: over a common length of twelve letters the
Stern--Brocot products for the patterns $(1,1),(2,2),(1,2)$ have traces
$322$, $198$ and $194$, so the alternating pattern grows strictly more slowly
than the Pell pattern. We therefore make no extremality claim from below.
\end{remark}

\subsection{Dictionary}

The correspondence between the word-combinatorial objects and the objects of
this paper is exact, and is summarised in \cref{tab:dictionary}. Its first row
is an identity rather than an analogy: by \cref{cor:writing-letter-counts} the
letter counts are read off the writing by a congruence, and the two orderings
of the writing are the two letter counts.

\begin{table}[H]
\centering
\caption{The word-combinatorial reading of the objects of this paper.}
\label{tab:dictionary}
\begin{tabular}{@{}ll@{}}
\toprule
Christoffel--Markoff side & This paper\\
\midrule
valid writing $m=p^2+q^2$ & seed $z=(u,v)$ with $\mathcal H(z)=u^2+v^2$\\
one step along the Pell arm & $\mathsf S(u,v)=(v,u+2v)$\\
that step, in the unit & multiplication by $\lambda$\\
Fibonacci arm & Fibonacci core $z_{q,\sigma}$\\
letter counts $(m_a,m_b)$ & coordinates of $\varphi^{2k}$ or $\lambda^{2k}$\\
cut values of the half directive word & the coefficients $F_r$, or $C_r$ and $P_r$\\
directive-word cost & the exponents $q$ and $n$\\
\bottomrule
\end{tabular}
\end{table}

The dictionary is already at work in this paper, without having been named.
The identity $P_M=F_{q+1}^2+F_{q+2}^2=F_{2q+3}$ of \cref{eq:jminus1} says, in
the left-hand column, that a Fibonacci-arm Markoff number, written as a sum
of two squares, is a Pell number; and it is settled there by Alekseyev's
Fibonacci--Pell intersection theorem \cite[Theorem~9]{Alekseyev2011}.

\subsection{Consequence for the Markoff-facing programme}

\Cref{thm:half-step-criterion,prop:no-half-step-34} close one route that
would otherwise look promising. Given the general orbit/gap equivalence of
\cref{thm:general-orbit-gap}, it is natural to try to prove the analogue of
\cref{thm:classification-intro} on another branch of the Markoff tree, by
transporting $\mathsf S$ there. That cannot be done: outside $m_0=1$ and
$m_0=2$ there is no integral half-step to transport, because
\cref{eq:half-step-criterion} fails, or, in the one remaining case, because
the branch has the wrong initial data. Any extension must therefore change
the mechanism, not merely its parameters.
